%=============================================================================
%  LNCS version - English
%  Text updated to match the corrected SBR version (main_corrigido.tex/pdf),
%  which incorporates the reviewers' comments: rewritten abstract, expanded
%  related work, revised adaptation section (six strategies under a common
%  protocol) and out-of-sample (two-fold CV) detection results.
%=============================================================================
\documentclass[runningheads]{llncs}

\usepackage[T1]{fontenc}
\usepackage[utf8]{inputenc}
\usepackage[english]{babel}
\usepackage{graphicx}
\usepackage{amsmath, amssymb}
\usepackage{booktabs}
\usepackage{subcaption}
\usepackage{multirow}

\graphicspath{{figs/}}

\emergencystretch=1.5em

\begin{document}

\title{Data Drift in a Seq2Seq Model for Robot Trajectory Prediction in the RoboCup SSL}
\titlerunning{Data Drift in Seq2Seq SSL Trajectory Prediction}

\author{Ivan G. Yamasaki\inst{1}\orcidID{0009-0000-4776-134X} \and
Marcos R. O. A. M\'aximo\inst{1}\orcidID{0000-0003-2944-4476} \and
Paulo M. Tasinaffo\inst{1}\orcidID{0000-0001-7069-8430}}
%
\authorrunning{I. Yamasaki et al.}
%
\institute{Aeronautics Institute of Technology, S\~ao Jos\'e dos Campos, Brazil}
%
\maketitle

\begin{abstract}
Trajectory prediction models are normally reported on a single fixed split of
the data, which says little about how long they stay useful. We take a
sequence-to-sequence (seq2seq) predictor for the RoboCup Small Size League
(SSL), trained on logs from the 2019 tournament, and evaluate it without any
retraining on the tournaments of 2021 to 2025. The point of comparison is a
baseline built on a Kalman filter for tracking and prediction, the classical
choice in the league. Accuracy does fall as the league moves away from the
training season. The Kalman-filter baseline falls faster, though, so the margin
in favor of the learned model widens instead of closing. Importance weighting
attributes 59--72\,\% of the extra error to a change in the input distribution
rather than in the input--output relation, and points to tail acceleration as
the main culprit; a classifier two-sample test agrees, season by season. We put
that diagnosis to work twice: it tells us how to weight the replay buffer of a
light fine-tuning, which then recovers accuracy with the least forgetting, and
it yields a monitoring signal that needs no model evaluation and stays below
0.13 false alarms per thousand trajectories on the drift-free training season.
\keywords{trajectory prediction \and data drift \and covariate shift \and drift detection \and RoboCup SSL \and mobile robots}
\end{abstract}

\section{Introduction}
\label{sec:intro}

Predicting where a robot will be a fraction of a second from now is a routine
requirement in systems with fast dynamics. In the RoboCup Small Size League
(SSL), where six to eleven robots per team move at several meters per second
under centralized control, such predictions feed planning, positioning and
real-time decision making. All matches used here come from Division~A, the top
division of the league, where the most mature teams play and where hardware and
playing styles change fastest.

The Kalman filter~\cite{kalman} is the traditional tool for this task, and it
is hard to beat on cost: deterministic, cheap and easy to tune. We call
\emph{Kalman-filter baseline} the constant-velocity predictor given by its
steady state, and use it throughout as the model-based reference. Its weakness
is the assumption behind it, namely that the dynamics are smooth enough for a
local linear model. Sequence-to-sequence (seq2seq) networks drop that
assumption and learn the patterns of future motion straight from logged
matches; Steuernagel and M\'aximo~\cite{seq2seq} showed that such a model beats
the kinematic reference in the SSL.

Being more accurate in one season, though, is not the same as still being
useful three seasons later. The league moves fast. Robots are rebuilt with
stronger motors, control stacks are rewritten, and playing styles change from
one tournament to the next. A network that learned the motion patterns of an
old season can decay silently under all of this, whereas a model-based
predictor has no learned parameters and might be expected to survive untouched.
That expectation is worth testing, and testing it is the data drift problem we
study.

Our contribution is threefold: we measure how a seq2seq trajectory predictor
for the SSL ages over seven years of official RoboCup logs, we identify the
mechanism behind the loss, and we turn that mechanism into adaptation and
monitoring procedures that can be deployed. Relative to~\cite{seq2seq}, the
question changes. It is no longer whether the learned model beats the kinematic
reference on a fixed split, but how far that advantage survives as the data
drift away from the training season. Four questions organize the study:

\begin{enumerate}
  \item how the error evolves across seasons, at two prediction horizons, and
        how that evolution compares to the Kalman-filter baseline;
  \item what fraction of the degradation is explained by covariate shift, and
        which kinematic dimension dominates it;
  \item whether light adaptation recovers accuracy without catastrophic
        forgetting;
  \item whether an online detector can flag the drift with a false-alarm rate
        low enough to trigger retraining.
\end{enumerate}

The paper is organized as follows. Section~\ref{sec:related} reviews related
work and Section~\ref{sec:method} describes the data, the model and the metric.
Section~\ref{sec:degradation} measures the degradation over time, and
Section~\ref{sec:iw} splits it into covariate shift and concept drift.
Sections~\ref{sec:adaptation} and~\ref{sec:detection} deal with the two
responses, adaptation and online detection. Section~\ref{sec:conclusion}
concludes.

\section{Related Work}
\label{sec:related}

\paragraph{Trajectory prediction.}
Most progress on learned trajectory predictors has come from the pedestrian and
driving domains. Social-LSTM~\cite{alahi2016} was among the first to model
interaction explicitly, pooling the hidden states of neighboring LSTMs so that
each agent is predicted with some awareness of the others.
Trajectron++~\cite{salzmann2020} replaced that fixed pooling with a dynamic
graph and added a latent variable, which produces several plausible futures
rather than one. AgentFormer~\cite{yuan2021} went further and let a transformer
attend over time and over agents at once, keeping social and temporal reasoning
together. Diffusion models~\cite{gu2022} have since been used to generate entire
trajectory distributions by progressive denoising. In the SSL, Steuernagel and
M\'aximo~\cite{seq2seq} applied an attention-based seq2seq model at the horizons
we also use and reported a clear gain over a kinematic predictor; that model is
our direct baseline. All of these evaluations rely on a fixed split, with
training and test data collected under the same conditions. Stability over time
is simply not part of the protocol.

\paragraph{Data drift and its detection.}
Gama et al.~\cite{gama2014} separate virtual drift, a change in $p(x)$, from
real drift, a change in $p(y|x)$. The difference has practical consequences:
virtual drift can in principle be corrected by reweighting data that is already
on disk, while real drift demands fresh labels. Applied work often skips this
step and reports only that accuracy dropped, although the remedy depends
entirely on which component moved. Detectors come in two families.
Error-based ones such as ADWIN~\cite{adwin} and Page--Hinkley~\cite{page_hinkley}
monitor the loss and therefore need labels; data-based ones such as
KSWIN~\cite{kswin} watch the inputs and do not. Lu et al.~\cite{lu2019} argue
that a detector only becomes trustworthy in production once its false-alarm rate
has been calibrated, and we adopt that requirement as a design constraint in
Section~\ref{sec:detection}. On the response side, adaptation is usually posed
as a stability--plasticity trade-off, with elastic weight
consolidation~\cite{kirkpatrick2017} as the standard regularizer against
catastrophic forgetting.

\paragraph{Covariate shift.}
When only $p(x)$ moves, the classical correction is importance weighting. Old
samples are reweighted by the ratio between the new and the old input density,
so that a quantity computed on old data behaves as if it had come from the new
one. Estimating the two densities separately is wasteful, and unstable already
in moderate dimension, which is why several methods estimate the ratio directly:
KLIEP~\cite{kliep} and least-squares importance fitting
(LSIF)~\cite{kanamori2009,sugiyama_book} are the best known.
RuLSIF~\cite{yamada2013} fits a relative ratio instead, accepting a small bias
in exchange for bounded weights and lower variance. The classifier two-sample
test~\cite{lopezpaz2017} offers an assumption-free alternative: train a
discriminator to tell the two samples apart and read the shift off its accuracy.
We use LSIF as the main estimator and the other two as robustness checks.

\section{Data and Model}
\label{sec:method}

\paragraph{Data.}
We use the official logs of the 2019--2025 RoboCup tournaments, with 2020
missing because the pandemic cancelled the competition altogether. Six
Division~A matches were taken from each tournament --- the same number every
year, so that no season weighs more than another --- for a total of 36 matches
recorded at 60\,Hz. The 2021 edition was played entirely in simulation, on
grSim, since the pandemic still ruled out an in-person tournament. We keep it in
the descriptive analyses, as it carries by far the largest distribution jump of
the period ($+3.05$\,mm of excess ADE), but part of that jump probably reflects
the gap between simulation and reality rather than organic drift. It is
therefore reported separately throughout, and left out of the adaptation
experiments of Section~\ref{sec:adaptation}: consolidating a model against an
anomalous simulated season would contaminate the old regime with a distribution
that no longer describes the league. Since every log comes from the same
division, the growing error cannot be blamed on a change in the competitive
level of the sample. The per-trajectory analyses draw 8\,000 windows per match,
which gives a temporally ordered stream of 288\,000 trajectories, 48\,000 per
year. All six matches of 2019 go into the baseline, and for good reason: a
sensitivity sweep puts $\mathrm{ADE}_{2019}$ anywhere between 4.53 and
6.20\,mm if a single reference match is used, enough to flip the sign of the
excess in four of the five remaining years.\footnote{Pipeline code and
processed data available from the authors upon request.}

\paragraph{Model and benchmark.}
The predictor follows the architecture of~\cite{seq2seq}: a 128-unit LSTM
encoder, the additive attention of Bahdanau et al.~\cite{bahdanau}, and an
autoregressive LSTM decoder without teacher forcing. Training uses
Adam~\cite{adam} with an initial learning rate of $10^{-3}$, batch size 1024
and early stopping. Every input step carries six normalized variables: the
position $x,y$ and the velocity $v_x,v_y$ of the robot in the horizontal plane
of the field, together with $\sin\psi$ and $\cos\psi$, where $\psi$ is the
heading angle about the vertical axis. We fit the normalization once, on the
2019 season, and apply it unchanged to every later tournament; otherwise the
later years would be measuring a change in preprocessing rather than in the
data. One detail of the reference model matters a great deal here. The target
lives in the space of normalized incremental velocities, and the future
position is rebuilt from the last observed state, so what the network really
predicts is accumulated dynamics. That makes it directly sensitive to the
velocity and acceleration distributions. The Kalman-filter baseline is the
constant-velocity predictor given by the steady state of a Kalman
filter~\cite{kalman}: it holds the last observed velocity and extrapolates the
position linearly. It is deterministic, has no tunable parameters and is
identical in every season, so there is no training distribution for it to drift
away from. Any error growth on its side measures how far the robots themselves
have moved away from constant-velocity behavior.

\paragraph{Metric.}
For $N$ trajectories and horizon $H$, the average displacement error is
\begin{equation}\label{eq:ade}
\mathrm{ADE} = \frac{1}{NH} \sum_{i=1}^{N}\sum_{t=1}^{H}
\left\lVert \hat{p}_{i,t} - p_{i,t} \right\rVert_2 ,
\end{equation}
with $\hat{p}_{i,t}$ and $p_{i,t}$ the predicted and the observed planar
position of trajectory $i$ at step $t$. The descriptive results cover two
horizons, $30 \rightarrow 15$ and $60 \rightarrow 30$, where $a \rightarrow b$
means observing $a$ past frames and predicting the next $b$ --- half a second
of history for a quarter of a second of future, at 60\,Hz. Everything that
follows on mechanism and detection uses $30 \rightarrow 15$.

\section{Temporal Degradation}
\label{sec:degradation}

\begin{figure}[t]
  \centering
  \includegraphics[width=\linewidth]{fig_degradacao_en}
  \caption{Temporal degradation (color = model, line style = horizon).
  (a) Mean ADE per year. (b) Percentage increase in ADE relative to 2019.
  (c) $\mathrm{p99}/\mathrm{p50}$ ratio of the per-trajectory ADE: the 2021
  jump is a tail effect ($4.36 \to 4.94$), whereas in 2022--2025 the ratio
  falls below baseline while the median rises by about 25\,\%.}
  \label{fig:degradation}
\end{figure}

Fig.~\ref{fig:degradation} summarizes three descriptive findings.

\paragraph{The seq2seq model degrades outside its training distribution.}
ADE grows in every season after 2019, but two effects have to be read apart.
The virtual 2021 season stands alone, with a jump of $+3.05$\,mm ($+67\,\%$)
for the reason given in Section~\ref{sec:method}. Set that season aside and
what remains over 2022--2025 is a slow upward trend against the same 2019
baseline: $+0.23$\,mm in 2022, $+0.73$\,mm in 2025 at the
$30 \rightarrow 15$ horizon, with 2024 almost back to zero. These are two
separate phenomena, not one stacked on the other: an abrupt regime change in a
simulated tournament, and a gradual drift in the in-person ones.
Both horizons show the same pattern.

\paragraph{The Kalman-filter baseline degrades more.}
Its relative error grows faster than the seq2seq model's in nearly every
season. The baseline has no learned parameters, so what deteriorates is not the
match to a training distribution but the match between the constant-velocity
assumption and the way the robots now move. This does not make the neural model
stable --- its absolute error grows too. It means the kinematic reference is
even more sensitive to the regime change, and the advantage of learning
therefore widens under drift.

\paragraph{The new regime enters through the tail, then spreads.}
The $\mathrm{p99}/\mathrm{p50}$ ratio rises only at the 2021 jump
($4.36 \to 4.94$), where the model fails disproportionately on a tail of
trajectories. In the seasons that follow, the ratio falls below baseline, to
$3.4$--$4.0$, because the median climbs by about 25\,\% up to 2025 while the
99th percentile stays flat. Behavior that used to be rare has become ordinary,
and the gradual degradation now sits in the body of the distribution.

\begin{figure}[t]
  \centering
  \includegraphics[width=\linewidth]{fig_drift_cinematico_en}
  \caption{Kinematic drift. (a) Kolmogorov--Smirnov statistic per year against
  the 2019 baseline (lines: annual mean; dots: individual matches); higher
  means a larger discrepancy between that year's distribution and 2019's.
  (b) Wasserstein distance indexed to the 2019 level, read the same way.
  (c) Mean ADE per match against the match's \texttt{accel\_p99}, with a
  regression line and a 95\,\% bootstrap band.}
  \label{fig:kinematic}
\end{figure}

Looking at the input side directly confirms the picture
(Fig.~\ref{fig:kinematic}). We use two distances between distributions, both
computed per year against the 2019 baseline: the Kolmogorov--Smirnov statistic,
which is the largest gap between the two cumulative distributions, and the
Wasserstein distance, which is the average displacement needed to turn one into
the other. Both panels are read the same way, a higher curve meaning a season
further from 2019. Velocity and acceleration show clear and persistent shifts;
turn rate moves much less. At match level, ADE grows almost monotonically with
the tail acceleration of the match (\texttt{accel\_p99}, Spearman
$\rho = 0.86$). The more aggressive the match, the worse the model does,
whatever season it belongs to.

\section{How Much of It Is Covariate Shift?}
\label{sec:iw}

The joint distribution factors as $p(x,y) = p(y|x)\,p(x)$~\cite{gama2014}.
Covariate shift moves only $p(x)$ and can be undone by reweighting old samples,
with no retraining at all; concept drift moves $p(y|x)$ and needs fresh labels.
Telling the two apart therefore decides whether the cheap fix is available.

LSIF~\cite{kliep,kanamori2009,sugiyama_book} estimates the density ratio
$\hat{w}(x) = p_{\text{target}}(x)/p_{2019}(x)$ directly as a combination of
Gaussian kernels, with coefficients $\boldsymbol\alpha$ obtained from
\begin{equation}\label{eq:lsif}
\min_{\boldsymbol\alpha \geq 0}\;
\tfrac{1}{2}\,\mathbb{E}_{p_{2019}}[\hat{w}(x)^2] -
\mathbb{E}_{p_{\text{target}}}[\hat{w}(x)] + \lambda\|\boldsymbol\alpha\|^2 ,
\end{equation}
with $x \in \mathbb{R}^{8}$ collecting the kinematic features of each observed
window: mean, p90 and p99 of speed and acceleration, plus mean and p90 of turn
rate, clipped at the 99.5th percentile of the pool and passed through
$\log(1+\cdot)$. Stability is monitored by the effective sample size
$\mathrm{ESS} = (\sum_i \hat{w}_i)^2/(n\sum_i \hat{w}_i^2)$, with the usual
$\mathrm{ESS} > 0.5$ rule of thumb~\cite{sugiyama_book}. The weighted error and
the recovered fraction of the excess are
\begin{equation}\label{eq:recovery}
\mathrm{ADE}_{\mathrm{IW}} = \frac{\sum_i \hat{w}(x_i)\,e_i}{\sum_i \hat{w}(x_i)},
\quad
\mathrm{rec.} = \frac{\mathrm{ADE}_y - \mathrm{ADE}_{\mathrm{IW}}}
                     {\mathrm{ADE}_y - \mathrm{ADE}_{2019}} ,
\end{equation}
where $e_i$ is the ADE of trajectory $i$ and $\mathrm{ADE}_y$ the
unweighted mean error of season $y$.
A recovery of 100\,\% would mean pure covariate shift, and 0\,\% pure concept
drift.

\begin{table}[t]
\centering
\caption{Importance-weighting decomposition (LSIF) with block bootstrap by
match ($n_{\text{boot}} = 2000$, 6 blocks per year). Baseline
$\mathrm{ADE}_{2019} = 4.53$\,mm; corr.\ $= \mathrm{ADE}_y -
\mathrm{ADE}_{\mathrm{IW}}$. 2024 has negative excess and is excluded from the
recovery column. All classifier tests give $p < 10^{-3}$ by permutation.
$^{\dagger}$2022 has a small excess (0.23\,mm), so the percentage should be
read with care.}
\label{tab:iw}
\small
\setlength{\tabcolsep}{3pt}
\begin{tabular}{@{}rrrlll@{}}
\toprule
\textbf{Year} & \textbf{ADE$_y$} & \textbf{ADE$_{\mathrm{IW}}$}
 & \textbf{rec.\ (\%) [95\,\% CI]} & \textbf{corr.\ (mm)} & \textbf{AUC} \\
\midrule
2021 & 7.58 & 5.70 & 61.6 $[43.7;\,79.3]$          & 1.88 $[1.33;\,2.42]$ & 0.754 \\
2022 & 4.75 & 4.62 & 59.0 $[9.2;\,115.2]^{\dagger}$ & 0.13 $[0.02;\,0.26]$ & 0.647 \\
2023 & 5.21 & 4.72 & 71.8 $[43.8;\,109.2]$          & 0.49 $[0.30;\,0.75]$ & 0.657 \\
2024 & 4.50 & 4.24 & ---                            & ---                  & 0.612 \\
2025 & 5.26 & 4.75 & 69.3 $[57.2;\,79.0]$           & 0.51 $[0.42;\,0.58]$ & 0.710 \\
\bottomrule
\end{tabular}
\end{table}

\paragraph{Robustness.}
Every number above depends on the weights behaving under the heavy tail of
\texttt{accel\_p99}, so we ran three checks. First, the weights stay bounded:
$\max_i \hat{w}_i \in [2.9;\,4.4]$ in every season, and the ESS of
$0.58$--$0.73$ comes from the near-zero mass given to the new regime, which is
the right mechanism for reconstructing 2019, not from a handful of extreme
weights (the top 1\,\% of trajectories carries only 3--4\,\% of the mass).
Second, capping the weights at $c \in \{1.5,\ldots,50\}$ leaves recovery
saturated from about $c = 3$ onward. Third, the RuLSIF relative-ratio transform
$\hat{w}/(\alpha\hat{w}+1-\alpha)$ with $\alpha = 0.1$ keeps recovery at
59--65\,\%. We also tried fitting RuLSIF directly and discarded it: in this
configuration it returns almost uniform weights ($\mathrm{ESS} \geq 0.94$) and
underestimates recovery at 9--23\,\%, which the two previous checks identify as
an artifact of the fit rather than a property of the estimand.

\paragraph{Reading.}
Table~\ref{tab:iw} puts LSIF recovery at 59--72\,\% in the years with a
positive excess, with confidence intervals that exclude zero. The absolute
correction supports a stronger statement, since its lower bound stays positive
even in 2022 ($\geq 0.02$\,mm): with only six bootstrap blocks, reweighting
still moves the error toward the 2019 level by an amount distinguishable from
zero. The classifier two-sample test agrees. It separates each season from 2019
with $\mathrm{AUC} \in [0.61;\,0.75]$ and ranks the seasons by severity in the
same order, without estimating any density ratio. What is left unrecovered,
30--40\,\%, is residual concept drift plus sampling noise.

\begin{table}[t]
\centering
\caption{Marginal recovery per feature (LSIF refitted on one feature at a
time), for the three seasons with a significant excess.}
\label{tab:iwfeat}
\small
\begin{tabular}{lcccc}
\toprule
\textbf{Feature} & \textbf{2021 (\%)} & \textbf{2023 (\%)} & \textbf{2025 (\%)}
 & \textbf{mean ESS} \\
\midrule
\texttt{accel\_p99}  & \textbf{57.1} & 52.6 & 51.7 & 0.55 \\
\texttt{accel\_p90}  & 54.0          & 51.3 & 51.1 & 0.57 \\
\texttt{accel\_mean} & 43.0          & 47.0 & 49.5 & 0.68 \\
\texttt{speed\_p99}  & 17.2          & 36.1 & 40.1 & 0.84 \\
\texttt{speed\_p90}  & 15.8          & 34.4 & 38.0 & 0.85 \\
\texttt{speed\_mean} &  7.9          & 24.7 & 27.4 & 0.91 \\
\texttt{turn\_p90}   &  3.9          &  0.2 &  0.0 & 0.97 \\
\texttt{turn\_mean}  &  5.3          &  0.9 &  0.0 & 0.96 \\
\bottomrule
\end{tabular}
\end{table}

Refitting LSIF on one feature at a time makes the mechanism concrete
(Table~\ref{tab:iwfeat}). Tail acceleration leads in every season, with
\texttt{accel\_p90} close behind; the speed features recover far less, and the
turn-rate features are essentially inert, below 5\,\% throughout. The ranking
barely changes from year to year. The league has adopted more aggressive
acceleration profiles, and that is the main vector of degradation. The ESS
column moves in the opposite direction, as expected: the features that shifted
the most are the ones whose reweighting costs the most effective sample.

\section{Adaptation}
\label{sec:adaptation}

Since the excess is mostly covariate, we compared four families of adaptation
strategies under a single conservative protocol: old regime $=$ 2019, new data
$=$ 2022--2025 (500 trajectories per match, 2021 excluded as an outlier),
frozen encoder, two trainable decoder layers, $\mathrm{lr} = 10^{-5}$ and five
epochs. Forgetting is measured by \texttt{cat\_ratio}, the fraction of old-regime
trajectories that get worse after the update. Table~\ref{tab:adapt} collects
the results.

\begin{table}[t]
\centering
\caption{Adaptation strategies under the common conservative protocol. The
unadapted model scores 4.584\,mm on the old regime and 4.965\,mm on the new
data. EWC reports the median of the $\lambda \in [50;\,10^7]$ plateau. The
last three rows are driven by the covariate diagnosis.}
\label{tab:adapt}
\small
\setlength{\tabcolsep}{4pt}
\begin{tabular}{@{}lccc@{}}
\toprule
\multirow{2}{*}{\textbf{Strategy}} & \multirow{2}{*}{\textbf{\texttt{cat\_ratio}} $\downarrow$}
 & \multicolumn{2}{c@{}}{\textbf{ADE (mm)}} \\
\cmidrule(l){3-4}
 & & \textbf{old} & \textbf{new} \\
\midrule
Conservative FT ($r=0$)     & 0.457          & 4.480          & 4.762 \\
EWC ($\lambda$ plateau)     & 0.457          & 4.481          & 4.817 \\
Uniform replay ($r=0.5$)    & 0.438          & 4.439          & \textbf{4.760} \\
\midrule
Importance-weighted replay  & \textbf{0.417} & \textbf{4.404} & 4.761 \\
Acceleration-selective FT   & 0.467          & 4.526          & 4.811 \\
Selective $+$ IW replay     & 0.449          & 4.441          & 4.800 \\
\bottomrule
\end{tabular}
\end{table}

Naive fine-tuning is destructive: without the conservative protocol,
\texttt{cat\_ratio} reaches 0.82 after a single epoch. Elastic weight
consolidation~\cite{kirkpatrick2017} is the standard remedy and adds
$\frac{\lambda}{2}\sum_i F_i(\theta_i - \theta^{*}_i)^2$, where $\theta_i$ is
the $i$-th parameter, $\theta^{*}_i$ its value before the update, and $F_i$ the
corresponding entry of the diagonal Fisher information, which says how much the
old task depends on that parameter. In our setting the Fisher term does
nothing. Sweeping $\lambda \in \{0, 50, \ldots, 10^7\}$ leaves
\texttt{cat\_ratio} pinned near 0.46 across eight orders of magnitude, because
the diagonal is numerically tiny ($F_i \approx 10^{-7}$), as one would expect
from an already converged decoder. The old regime is protected by the small
learning rate and the frozen encoder, not by the penalty.

Mixing a replay buffer of original data into the update, at a replay-to-new
ratio $r = 0.5$, improves all three metrics at once: \texttt{cat\_ratio} drops
below the no-replay control and the ADE on new data is the lowest in the table.
The gain is modest, since the conservative schedule already limits forgetting,
but it is consistent and costs essentially nothing.

\begin{figure}[t]
  \centering
  \includegraphics[width=\linewidth]{fig_finetune_en}
  \caption{Adaptation strategies under the same protocol (blue: variants driven
  by the covariate/acceleration diagnosis). (a) Forgetting, with the error bar
  showing the EWC $\lambda$ plateau and the dotted line the no-replay control.
  (b) ADE after adaptation on both regimes, with the unadapted model as
  reference (dashed lines).}
  \label{fig:finetune}
\end{figure}

Steering the adaptation by the covariate diagnosis gives a more interesting
split (Fig.~\ref{fig:finetune}). Sampling the replay buffer in proportion to
$\hat{w}(x)$ helps: it reaches the lowest \texttt{cat\_ratio} (0.417) and the
lowest old-regime ADE in the table, albeit by a small margin and on a single
seed. Hard selection, on the other hand, hurts. Fine-tuning restricted to the
top 50\,\% of new data by tail acceleration --- exactly the region where $p(x)$
moved --- gives the worst forgetting of all six rows (0.467), worse even than
the control. The tail is the subset farthest from the 2019 regime, so training
on it alone displaces the weights more. The unchanged body of the distribution,
which hard selection throws away, was acting as an implicit anchor; reweighting
keeps that anchor and merely emphasizes the relevant tail. The diagnosis is thus
worth using for monitoring, for explaining the excess and for weighting replay,
but not for deciding which data to discard.

\section{Online Detection}
\label{sec:detection}

The operational question is whether drift can be flagged automatically on the
288\,k-trajectory stream. We judge a detector by three quantities~\cite{lu2019}:
$\mathrm{FAR}_{2019}$, the alarms per thousand trajectories in the training
year, where no drift exists and every alarm is by definition false
(admissibility ceiling 0.20/1k); coverage, the number of 2021--2025 seasons with
at least one alarm; and latency after each year boundary. To avoid tuning and
reporting on the same data, calibration runs under a two-fold cross-validation
by matches: half the matches of each year select the ADWIN~\cite{adwin}
configuration out of a 96-point grid, and the untouched half is evaluated once.
Every figure below is out of sample.

The detector is only viable at all thanks to a local robust normalization of the
per-trajectory signal, by median and MAD over a 200-trajectory window. On the
raw signal, no configuration in the grid is admissible in either fold, and the
21 conventional smoothers we tried (moving averages, EWMA, Holt,
Savitzky--Golay) all raise FAR by injecting the very autocorrelation that ADWIN
assumes away. On the normalized ADE the winning family is stable across folds
($\delta = 10^{-4}$, windows 1000/5000): $\mathrm{FAR}_{2019}$ of 0.042/1k with
5/5 coverage in fold~A, 0.125/1k with 4/5 in fold~B, at a latency of 0.2 to 1.5
matches. The holdout does more than add confidence here, it changes what one
would deploy: the configuration picked by in-sample selection keeps
$\mathrm{FAR} \approx 0$ and then covers only 2--3 of the 5 seasons out of
sample.

The diagnosis of Section~\ref{sec:iw} also buys a second monitoring channel.
Since the shift is dominated by tail acceleration, the features of the observed
window already carry the drift signal, with no need to run the model or to wait
for the 15 future steps that ADE requires. Applying the same protocol to four
kinematic signals reproduces the ranking of the per-feature decomposition, with
\texttt{accel\_p95} in front ($\mathrm{FAR} = 0.25$--$0.29$/1k, 5/5 coverage in
both folds).\footnote{p95 rather than the dominant p99, because the observed
window holds only 29 acceleration samples.} No single feature is admissible on
its own, because kinematics vary naturally within 2019 as well. Still, the
acceleration channel fires before the error channel in 8 of 10 year-folds, which
is enough for a watch layer. At the other extreme, requiring both alarms within
a coincidence window of 200 trajectories drives observed false alarms to zero in
both folds, at the cost of coverage (2--3 of 5).

\begin{table}[t]
\centering
\caption{Layered detection architecture, with metrics on the test split of
each fold (A/B). All signals use robust normalization with $w = 200$; each
detector is configured per fold during calibration.}
\label{tab:detection}
\small
\setlength{\tabcolsep}{3pt}
\begin{tabular}{@{}llccl@{}}
\toprule
\textbf{Layer} & \textbf{Signal} & \textbf{FAR/1k (A/B)} & \textbf{Cov.}
 & \textbf{Action} \\
\midrule
Primary      & seq2seq ADE          & 0.042 / 0.125 & 5/5, 4/5 & retrain \\
Watch        & \texttt{accel\_p95}  & 0.292 / 0.250 & 5/5, 5/5 & monitor \\
Confirmation & AND(ADE, accel)      & 0.000 / 0.000 & 3/5, 2/5 & confirm \\
\midrule
Corroborator & Kalman-filter ADE  & 0.167 / 0.083 & 5/5, 5/5 & redundancy \\
\bottomrule
\end{tabular}
\end{table}

Table~\ref{tab:detection} summarizes the resulting three-layer architecture.
Two caveats and one negative result should be recorded. Per-trajectory FAR
assumes roughly independent trajectories within 2019, and the serial
correlation is certainly there ($\rho_1 = 0.54$); a permutation test suggests
the violation is conservative, since shuffling trajectories within matches
drives false alarms to zero in all 200 replicates. Among the alternatives,
Page--Hinkley~\cite{page_hinkley} aggregated over windows achieves
$\mathrm{FAR} = 0$ but catches only the abrupt 2021 jump, while
KSWIN~\cite{kswin} is inadmissible in every configuration
($\mathrm{FAR} \geq 0.438$/1k). The negative result is structural. Combining
detectors that run at different time scales simply does not work here: the
closest pair of alarms is 214 trajectories apart, beyond any coincidence window
one would accept. Aligning the scales first removes the problem, and that is
how the confirmation layer above was built.

\section{Conclusions and Future Work}
\label{sec:conclusion}

A seq2seq trajectory predictor for the SSL loses accuracy as the league moves
away from its training season. The loss is gradual, it is measurable, and most
of it traces back to one specific change in the input distribution. Excess ADE
trends upward over 2022--2025, separately from the one-off jump of the virtual
2021 edition, and the Kalman-filter baseline degrades faster still. Importance
weighting recovers between 59 and 72\,\% of that excess, an estimate that
survives weight clipping and the relative-ratio transform and is corroborated
by an independent classifier test. Tail acceleration is the dominant vector.

On the response side, light adaptation keeps forgetting under control
(\texttt{cat\_ratio} from 0.82 to 0.46), EWC turns out to be inert over eight
orders of magnitude of $\lambda$, and a replay buffer improves forgetting and
adaptation at once, best of all when sampled by importance weights. Hard
selection by the drifting tail backfires. On the operational side, ADWIN over a
robustly normalized error stream reaches 0.04--0.13 false alarms per thousand
trajectories out of sample, covering 9 of 10 year-folds, and the same diagnosis
yields a model-free acceleration channel that anticipates the error alarm in 8
of them.

Several limitations bound these conclusions. Only Division~A appears in the
logs. Six matches per season keep the block confidence intervals wide, and the
Wilson intervals of the holdout cross the admissibility ceiling, so a claim of
admissibility at 95\,\% confidence would need more 2019 matches. ADE is
single-mode and would have to give way to
$\text{min-ADE}_K$~\cite{salzmann2020,yuan2021} for multimodal predictors.
Robust normalization assumes gradual drift and would absorb part of an abrupt
change. The residual concept-drift fraction is not isolated per feature, and it
may hide covariate shift in dimensions we did not measure. The 2021 season,
finally, was played in simulation, so its error jump mixes the gap between
simulation and reality with seasonal drift; the conclusions about the gradual
trend rest on 2022--2025 alone. The natural next step is a deeper adaptation,
retraining the encoder itself with an importance-weighted replay buffer.

\section*{Acknowledgements}
Ivan Yamasaki acknowledges the support of ITAEx and Nubank, provided within the framework of the ITA Undergraduate Research Program in Artificial Intelligence (2025--2026).

\begin{thebibliography}{9}

\bibitem{seq2seq}
Steuernagel, L., M\'aximo, M. R. O. A.: Trajectory Prediction for SSL Robots Using Seq2seq Neural Networks. In: Eguchi, A., Lau, N., Paetzel-Pr\"usmann, M., Wanichanon, T. (eds.) RoboCup 2022: Robot World Cup XXV. LNCS, vol. 13561, pp. 27--38. Springer, Cham (2023). doi:10.1007/978-3-031-28469-4\_4

\bibitem{bahdanau}
Bahdanau, D., Cho, K., Bengio, Y.: Neural Machine Translation by Jointly Learning to Align and Translate. In: 3rd International Conference on Learning Representations (ICLR) (2015)

\bibitem{kalman}
Kalman, R. E.: A New Approach to Linear Filtering and Prediction Problems. Journal of Basic Engineering \textbf{82}(1), 35--45 (1960)

\bibitem{adam}
Kingma, D. P., Ba, J.: Adam: A Method for Stochastic Optimization. In: 3rd International Conference on Learning Representations (ICLR) (2015)

\bibitem{kliep}
Sugiyama, M., Nakajima, S., Kashima, H., von B\"unau, P., Kawanabe, M.: Direct importance estimation with model selection and its application to covariate shift adaptation. In: Advances in Neural Information Processing Systems (NIPS) (2008)

\bibitem{sugiyama_book}
Sugiyama, M., Kawanabe, M.: Machine Learning in Non-Stationary Environments: Introduction to Covariate Shift Adaptation. MIT Press (2012)

\bibitem{kanamori2009}
Kanamori, T., Hido, S., Sugiyama, M.: A least-squares approach to direct importance estimation. Journal of Machine Learning Research \textbf{10}, 1391--1445 (2009)

\bibitem{yamada2013}
Yamada, M., Suzuki, T., Kanamori, T., Hachiya, H., Sugiyama, M.: Relative density-ratio estimation for robust distribution comparison. Neural Computation \textbf{25}(5), 1324--1370 (2013). doi:10.1162/NECO\_a\_00442

\bibitem{lopezpaz2017}
Lopez-Paz, D., Oquab, M.: Revisiting Classifier Two-Sample Tests. In: 5th International Conference on Learning Representations (ICLR) (2017). arXiv:1610.06545

\bibitem{gama2014}
Gama, J., \v{Z}liobait\.{e}, I., Bifet, A., Pechenizkiy, M., Bouchachia, A.: A survey on concept drift adaptation. ACM Computing Surveys \textbf{46}(4), 1--37 (2014). doi:10.1145/2523813

\bibitem{lu2019}
Lu, J., Liu, A., Dong, F., Gu, F., Gama, J., Zhang, G.: Learning under concept drift: a review. IEEE Transactions on Knowledge and Data Engineering \textbf{31}(12), 2346--2363 (2019). doi:10.1109/TKDE.2018.2876857

\bibitem{adwin}
Bifet, A., Gavald\`a, R.: Learning from time-changing data with adaptive windowing. In: Proceedings of the 2007 SIAM International Conference on Data Mining (SDM), pp. 443--448 (2007). doi:10.1137/1.9781611972771.42

\bibitem{page_hinkley}
Page, E. S.: Continuous Inspection Schemes. Biometrika \textbf{41}(1--2), 100--115 (1954). doi:10.1093/biomet/41.1-2.100

\bibitem{kswin}
Raab, C., Heusinger, M., Schleif, F.-M.: Reactive Soft Prototype Computing for Concept Drift Streams. Neurocomputing \textbf{416}, 340--351 (2020). doi:10.1016/j.neucom.2020.01.119

\bibitem{kirkpatrick2017}
Kirkpatrick, J., Pascanu, R., Rabinowitz, N., et~al.: Overcoming catastrophic forgetting in neural networks. Proceedings of the National Academy of Sciences \textbf{114}(13), 3521--3526 (2017). doi:10.1073/pnas.1611835114

\bibitem{alahi2016}
Alahi, A., Goel, K., Ramanathan, V., Robicquet, A., Fei-Fei, L., Savarese, S.: Social LSTM: Human trajectory prediction in crowded spaces. In: IEEE Conference on Computer Vision and Pattern Recognition (CVPR), pp. 961--971 (2016). doi:10.1109/CVPR.2016.110

\bibitem{salzmann2020}
Salzmann, T., Ivanovic, B., Chakravarty, P., Pavone, M.: Trajectron++: Dynamically-feasible trajectory forecasting with heterogeneous data. In: European Conference on Computer Vision (ECCV) (2020). arXiv:2001.03093

\bibitem{yuan2021}
Yuan, Y., Weng, X., Ou, Y., Kitani, K.: AgentFormer: Agent-aware transformers for socio-temporal multi-agent forecasting. In: IEEE/CVF International Conference on Computer Vision (ICCV) (2021). arXiv:2103.14023

\bibitem{gu2022}
Gu, T., Chen, G., Li, J., Lin, C., Rao, Y., Zhou, J., Lu, J.: Stochastic Trajectory Prediction via Motion Indeterminacy Diffusion. In: IEEE/CVF Conference on Computer Vision and Pattern Recognition (CVPR) (2022). arXiv:2203.13777

\end{thebibliography}

\end{document}
